\documentclass[11pt]{article}
\usepackage[margin=0.85in]{geometry}
\usepackage{amsmath,amssymb,mathtools,booktabs,longtable,array,enumitem,hyperref,bm,microtype}
\hypersetup{colorlinks=true,linkcolor=blue,citecolor=blue,urlcolor=blue}
\newcommand{\R}{\mathbb{R}}
\newcommand{\argmax}{\operatorname*{arg\,max}}
\newcommand{\lexmin}{\operatorname*{lex\,min}}
\newcommand{\lexmax}{\operatorname*{lex\,max}}
\title{\textbf{OpenDSSSim PESGM--2027}\\P2 Reconfiguration/Restoration Control Science\\\textbf{FROZEN Mathematical and Software Contract}}
\author{Controller-science track}
\date{Freeze version v01R1 --- 27 October 2026}
\begin{document}
\maketitle
\begin{center}
\begin{minipage}{0.94\linewidth}
\textbf{Status: FROZEN FOR IMPLEMENTATION.}
This document consolidates and supersedes the P2 Blocks 1--11 ratification discussion.
It freezes the common P2 physical problem, graph/reference controller, LinDistFlow
MILP MPC controller, masked-DQN RL controller, fair-information contract,
controller-side orchestration boundary, and reproducibility contract. It is a
science/design contract, not controller implementation code.
\end{minipage}
\end{center}
\tableofcontents
\newpage

\section{Scope, contribution boundary, and provenance}
The PESGM--2027 P2 study compares three controller paradigms on the same post-fault
distribution-system restoration problem: graph/reference restoration, mixed-integer
model predictive control (MPC), and reinforcement learning (RL).

The controller algorithms are not claimed as algorithmic novelties. The generalized
claim is: \emph{a generalized controller implementation can instantiate for any
compatible OpenDSS circuit satisfying its declared capability requirements}. It is
not claimed that every controller works on every OpenDSS circuit.

Primary source anchors are Vedullapalli, Hadidi, and Bozeman~\cite{Vedullapalli2020}
for graph restoration after multiple faults, Liu, Ding, and Zhao~\cite{Liu2020} for
dynamic/receding-horizon restoration using MPC, and Oh, Yoon, and Kim~\cite{Oh2020}
for online network-topology selection using deep Q-learning.

\textbf{REUSE} denotes source-backed ideas retained in the design.
\textbf{ADAPT} denotes deliberate changes needed so all three paradigms solve the
same P2 problem. \textbf{NEW software abstraction} denotes software infrastructure
without an algorithmic novelty claim.

\section{Block 1: common P2 physical restoration problem}
P2 is frozen as a post-fault, multi-fault, radial distribution-system
service-restoration problem. Faulted components are known at the restoration
boundary and unavailable. Healthy but de-energized sections are restoration targets.

The primary benchmark includes multiple simultaneous faults, post-isolation
restoration, sectionalizing/tie-switch control, radial source-connected operation,
priority-weighted service restoration, explicit controllable load curtailment where
declared, and voltage/thermal operational constraints.

The primary benchmark excludes crew routing, repair scheduling, protection redesign,
fault-detection optimization, grid-forming DER islands, frequency restoration,
storage dispatch, and DER dispatch. Utility/source-side restoration is primary.

The common objective philosophy is
\[
\boxed{\text{restore the most important healthy load safely, then minimize switching}.}
\]

\section{Block 2: graph, health, switching, source, and energization}
Represent the feeder by
\[
\boxed{\mathcal G=(\mathcal N,\mathcal E)}
\]
and partition edges by controllability:
\[
\boxed{\mathcal E=\mathcal E^{fix}\,\dot\cup\,\mathcal E^{sw}.}
\]
This distinction says nothing about health.

For every physical branch,
\[
\boxed{h_{e,t}\in\{0,1\}}
\]
with $h_{e,t}=1$ healthy/available and $h_{e,t}=0$ faulted/unavailable.
Health is exogenous.

For $e\in\mathcal E^{sw}$,
\[
\boxed{z_{e,t}\in\{0,1\},\qquad 0=\text{open},\quad 1=\text{closed}.}
\]

Define branch connectivity status for all edges:
\[
\boxed{
\delta_{e,t}=
\begin{cases}
h_{e,t}, & e\in\mathcal E^{fix},\\
h_{e,t}z_{e,t}, & e\in\mathcal E^{sw}.
\end{cases}}
\]

Let $\mathcal S\subseteq\mathcal N$ denote source buses and
$\sigma_{s,t}\in\{0,1\}$ their availability. Define
\[
\mathcal S_t^{avail}=\{s\in\mathcal S:\sigma_{s,t}=1\}.
\]
For every bus,
\[
\boxed{y_{i,t}\in\{0,1\}}
\]
denotes energization, with the semantic condition
\[
\boxed{
y_{i,t}=1
\iff
\exists s\in\mathcal S_t^{avail}\text{ and a healthy closed path }s\leadsto i.
}
\]

Healthy does not imply energized, and faulted is not synonymous with de-energized.
Healthy de-energized zones are the primary restoration target.

For every load define a physical served fraction
\[
\alpha_{\ell,t}\in[0,1],\qquad
\boxed{0\le\alpha_{\ell,t}\le y_{i(\ell),t}.}
\]

Controller-side problem construction partitions
\[
\boxed{\mathcal L=\mathcal L^{curt}\,\dot\cup\,\mathcal L^{fixed}.}
\]
Only $\mathcal L^{curt}$ appears as explicit continuous curtailment action channels.

\section{Block 3: restoration value, load priority, and common metrics}
Each load has service-priority weight
\[
\boxed{w_\ell\in\{1,2,3,4\}}
\]
from normal to very-high criticality. Exact assignments are experiment configuration.
Electrical distance is not embedded in $w_\ell$.

Priority-weighted restored and unserved load are
\[
R_t=\sum_\ell w_\ell\alpha_{\ell,t}P_{\ell,t}^L,
\qquad
U_t=\sum_\ell w_\ell(1-\alpha_{\ell,t})P_{\ell,t}^L.
\]

Define
\[
\boxed{
\rho_t^{restore}
=
\frac{\sum_\ell w_\ell\alpha_{\ell,t}P_{\ell,t}^{L}}
{\sum_\ell w_\ell P_{\ell,t}^{L}}
}
\]
and ordinary served-load fraction
\[
\boxed{
\rho_t^{kW}
=
\frac{\sum_\ell \alpha_{\ell,t}P_{\ell,t}^{L}}
{\sum_\ell P_{\ell,t}^{L}}.
}
\]

Over a horizon,
\[
\boxed{
E^{WU}=\sum_t\Delta t\sum_\ell w_\ell(1-\alpha_{\ell,t})P_{\ell,t}^{L}}
\]
and
\[
\boxed{
E^{U}=\sum_t\Delta t\sum_\ell(1-\alpha_{\ell,t})P_{\ell,t}^{L}.}
\]

Switching effort is
\[
\boxed{
J_{\mathrm{switch}}
=
\sum_t\sum_{e\in\mathcal E^{sw}}|z_{e,t}-z_{e,t-1}|.}
\]

The common hierarchy is service restoration before switching minimization.
Voltage and thermal requirements are operational constraints.

\section{Block 4: common action semantics and time}
At each P2 timestep the controller provides all P2 control channels:
\[
\boxed{
a_t^{P2,req}
=
\left(
\mathbf z_t^{req},
\boldsymbol\alpha_{\mathcal L^{curt},t}^{req}
\right).}
\]
Here $\mathbf z_t^{req}\in\{0,1\}^{N_{sw}}$ is the complete switch-state vector and
$\boldsymbol\alpha_{\mathcal L^{curt},t}^{req}\in[0,1]^{N_{curt}}$ contains only
declared curtailable-load channels. There is no requirement that $N_{sw}=N_{curt}$.

Multiple switch states may change between $t-1$ and $t$. There are no hidden
simulator-generated switching microsteps and no simulator-side break-before-make
sequence. Topology-transition logic belongs to the controller.

The common plant step is
\[
x_t\rightarrow o_t\rightarrow \text{controller}\rightarrow a_t^{req}
\rightarrow\text{admissibility}\rightarrow\text{apply admitted commands}
\rightarrow\text{one plant solve}\rightarrow x_{t+1}.
\]

Use a conventional fixed P2 timestep
\[
\boxed{\Delta t_{P2}\in\{5,15,30,60\}\ {\rm min}}
\]
with the exact value deferred to experiment freeze.

\section{Block 5: structural and operational feasibility}
Faulted switchable edges cannot be used:
\[
\boxed{z_{e,t}\le h_{e,t},\qquad e\in\mathcal E^{sw}.}
\]

Source connectivity and radiality are controller responsibilities. The energized
network must be source-connected and acyclic. Edge count alone is insufficient;
connectivity must also be enforced.

The authoritative plant is the nonlinear unbalanced AltDSS/OpenDSS solution.
Realized operational conditions are evaluated against declared limits such as
\[
\boxed{V_i^{min}\le |V_{i,\phi,t}|\le V_i^{max}}
\]
and
\[
\boxed{|I_{e,\phi,t}|\le I_{e,\phi}^{max}}
\]
or the corresponding native thermal quantity.

Power-flow convergence $c_t^{PF}\in\{0,1\}$ is recorded as a diagnostic, but is not
by itself the scientific definition of action success/failure.

OpenDSSSim admits only commands that satisfy the basic plant/action contract and
then reports what occurred. It does not restore radiality, select a replacement
topology, change $\boldsymbol\alpha$, optimize a substitute action, or hide violations.

\textbf{Slacks are permitted only on operational constraints.}
Fault availability, action validity, source connectivity, and radiality remain hard.

\section{Block 6: graph/reference restoration controller}
Denote the deterministic controller by $\pi^{GR}$.

At time $t$, controller-side orchestration provides current plant state/health plus
static circuit information:
\[
\mathcal I_t^{GR}
=
\{\mathcal G,\mathbf h_t,\mathbf z_{t-1}^{real},\boldsymbol\sigma_t,
\mathbf P_t^L,\mathbf Q_t^L,\mathbf w,\text{declared limits}\}.
\]

Construct the healthy graph
\[
\mathcal E_t^H=\{e\in\mathcal E:h_{e,t}=1\},\qquad
\mathcal G_t^H=(\mathcal N,\mathcal E_t^H),
\]
identify the current source-connected energized component(s), and form the healthy
de-energized components $\mathcal C_t^{off}=\{C_1,\dots,C_K\}$.

For component $C_k$ define restoration value
\[
\boxed{
R_k=\sum_{\ell:i(\ell)\in C_k}w_\ell P_{\ell,t}^L.}
\]

Following the graph-restoration philosophy of~\cite{Vedullapalli2020}, enumerate
healthy simple paths between healthy inactive sections and energized/source-connected
sections. Convert each candidate to a complete switch vector
\[
\mathbf z^{(c)}\in\{0,1\}^{N_{sw}}.
\]

Reject controller-side candidates that use unavailable edges, fail required source
connectivity, or create energized cycles.

For candidate $c$, compute graph-level restored value
\[
R^{(c)}=\sum_\ell w_\ell\hat y_\ell^{(c)}P_{\ell,t}^{L}
\]
and switching effort
\[
N_{\rm sw}^{(c)}
=
\sum_{e\in\mathcal E^{sw}}|z_e^{(c)}-z_{e,t-1}^{real}|.
\]
Rank
\[
\boxed{
\lexmax_c\left(R^{(c)},-N_{\rm sw}^{(c)}\right).}
\]

The controller may use a non-mutating controller-side current-state evaluator
\[
\boxed{
\mathcal E_{\rm eval}^{ctrl}
(x_t,\mathbf z^{(c)},\boldsymbol\alpha^{(c)})}
\]
to estimate voltage/loading and violations. This is not an OpenDSSSim service and
does not mutate the plant. The initial implementation may use a simpler
current-state multiphase LinDistFlow evaluation.

Among otherwise tied candidates, voltage quality may be used as a tertiary tie-break.

For the selected topology, the deterministic priority-aware curtailment decoder
starts with full service for reconnectable curtailable loads. If operational
violations are predicted, it curtails lower $w_\ell$ first; within equal priority
it chooses reductions that best relieve the active violation under the declared
heuristic, and stops at the minimum required heuristic curtailment.

The final output is
\[
\boxed{
a_t^{GR}
=
(\mathbf z_t^{req},\boldsymbol\alpha_{\mathcal L^{curt},t}^{req}).}
\]

\textbf{REUSE}: graph representation, known-fault isolation, active/inactive
subgraphs, path search, minimum-switch preference, operational checking,
priority-aware curtailment, and feasible-topology voltage tie-break
from~\cite{Vedullapalli2020}.

\textbf{ADAPT}: common lexicographic restoration metric, full-vector action,
controller-side radiality, utility-source primary benchmark, no crew scheduling,
no distance term in priority, explicit curtailment action, closed-loop recomputation.

\textbf{NEW software abstraction}: typed controller context, circuit/model compiler,
candidate evaluator, diagnostics, and provenance.

\section{Block 7: LinDistFlow MILP MPC controller}
Denote the controller by $\pi_{P2}^{MPC}$.

At time $t$, MPC optimizes $k=0,\dots,H-1$ and applies only the first decision:
\[
\boxed{
a_t^{MPC}
=
\left(
\mathbf z_{t|t}^{\star},
\boldsymbol\alpha_{\mathcal L^{curt},t|t}^{\star}
\right).}
\]
At the next time step it observes the authoritative plant and resolves.

The common load forecast is
\[
\hat{\mathbf D}_t
=
\{\hat P_{\ell,t+k|t}^{L},\hat Q_{\ell,t+k|t}^{L}\}_{k=0}^{H-1}.
\]
Observed faults are treated as persistent across the prediction horizon unless an
external scenario explicitly changes them. No future-fault oracle is used.

Decision variables include
\[
z_{e,k}\in\{0,1\},\quad
y_{i,k}\in\{0,1\},\quad
\alpha_{\ell,k}\in[0,1],
\]
phase-wise $P_{ij,k}^{\phi}$, $Q_{ij,k}^{\phi}$, squared voltage
$v_{i,k}^{\phi}=|V_{i,k}^{\phi}|^2$, energized-edge binary $g_{e,k}$, and
fictitious connectivity flow $F_{ij,k}\ge0$.

The variable $y_{i,k}$ is an auxiliary energization decision variable. It does not
directly actuate bus $i$; it is constrained so that it represents source connectivity
under the chosen topology.

For all physical edges,
\[
\boxed{
\delta_{e,k}=
\begin{cases}
h_{e,t}, & e\in\mathcal E^{fix},\\
h_{e,t}z_{e,k}, & e\in\mathcal E^{sw}.
\end{cases}}
\]

Distinguish physical closed/healthy status from participation in the energized tree
using
\[
g_{e,k}\in\{0,1\},\qquad e\in\mathcal E.
\]
For edge $e=(i,j)$,
\[
g_{e,k}\le \delta_{e,k},\qquad
g_{e,k}\le y_{i,k},\qquad
g_{e,k}\le y_{j,k},
\]
and
\[
g_{e,k}\ge \delta_{e,k}+y_{i,k}+y_{j,k}-2.
\]

For controllably curtailable load,
\[
\boxed{0\le\alpha_{\ell,k}\le y_{i(\ell),k}.}
\]
Use proportional $P/Q$ service to preserve power factor:
\[
p_{\ell,k}^{served}=\alpha_{\ell,k}\hat P_{\ell,k}^{L},\qquad
q_{\ell,k}^{served}=\alpha_{\ell,k}\hat Q_{\ell,k}^{L}.
\]

For the primary single-source formulation, set $y_{s,k}=1$ when the source is
available. Fictitious flow proves source connectivity:
\[
\boxed{
\sum_{j:(j,i)}F_{ji,k}-\sum_{j:(i,j)}F_{ij,k}=y_{i,k},
\qquad i\neq s,}
\]
\[
\boxed{
\sum_jF_{sj,k}-\sum_jF_{js,k}=\sum_{i\neq s}y_{i,k},}
\]
and
\[
\boxed{
0\le F_{ij,k}\le M_F g_{e(i,j),k}.}
\]

Connectivity flow alone does not prohibit loops, so enforce energized-tree
cardinality:
\[
\boxed{
\sum_{e\in\mathcal E}g_{e,k}
=
\sum_{i\in\mathcal N}y_{i,k}-1.}
\]
Together with the connectivity equations, this yields a source-connected acyclic
energized topology for the primary single-source case. A future generalized
multi-source implementation may use a super-root/forest extension.

The MPC prediction model is frozen as
\[
\boxed{\text{multiphase unbalanced LinDistFlow}.}
\]

For intuition, scalar exact radial DistFlow contains
\[
P_{ij}=p_j+\sum_{m:j\to m}P_{jm}+r_{ij}\ell_{ij},
\]
\[
Q_{ij}=q_j+\sum_{m:j\to m}Q_{jm}+x_{ij}\ell_{ij},
\]
\[
v_j=v_i-2(r_{ij}P_{ij}+x_{ij}Q_{ij})
+(r_{ij}^2+x_{ij}^2)\ell_{ij},
\]
\[
v_i\ell_{ij}=P_{ij}^2+Q_{ij}^2.
\]

LinDistFlow neglects the quadratic loss terms:
\[
P_{ij}\approx p_j+\sum_{m:j\to m}P_{jm},
\]
\[
Q_{ij}\approx q_j+\sum_{m:j\to m}Q_{jm},
\]
\[
v_j\approx v_i-2(r_{ij}P_{ij}+x_{ij}Q_{ij}).
\]

For the unbalanced three-phase controller model, retain cross-phase coupling:
\[
\boxed{
\mathbf v_{j,k}
=
\mathbf v_{i,k}
-
\mathbf R_{ij}\mathbf P_{ij,k}
-
\mathbf X_{ij}\mathbf Q_{ij,k}.}
\]
The matrices $\mathbf R_{ij}$ and $\mathbf X_{ij}$ are compiled controller-side
from the static OpenDSS circuit data and the adopted linearization convention.

Apply the voltage-drop relation only on energized tree edges:
\[
\boxed{
-M_V(1-g_{e,k})
\le
\mathbf v_{j,k}-\mathbf v_{i,k}
+\mathbf R_{ij}\mathbf P_{ij,k}
+\mathbf X_{ij}\mathbf Q_{ij,k}
\le
M_V(1-g_{e,k}).}
\]

Phase-wise power balance is
\[
\boxed{
\sum_{i:i\to j}P_{ij,k}^{\phi}
-
\sum_{m:j\to m}P_{jm,k}^{\phi}
=
\sum_{\ell\in\mathcal L_j}p_{\ell,k}^{served,\phi}}
\]
and similarly
\[
\boxed{
\sum_{i:i\to j}Q_{ij,k}^{\phi}
-
\sum_{m:j\to m}Q_{jm,k}^{\phi}
=
\sum_{\ell\in\mathcal L_j}q_{\ell,k}^{served,\phi}.}
\]
Implementation must use one consistent sign/orientation convention.

Operational voltage constraints may use slacks:
\[
\boxed{
(V_i^{min})^2y_{i,k}-\xi_{i,\phi,k}^{V,-}
\le
v_{i,\phi,k}
\le
(V_i^{max})^2y_{i,k}+\xi_{i,\phi,k}^{V,+}}
\]
with $\xi_{i,\phi,k}^{V,\pm}\ge0$.

Thermal limits use a conservative eight-sided polygon so the controller remains MILP.
Let
\[
\theta_m=\frac{(2m+1)\pi}{8},\qquad m=0,\dots,7.
\]
For each facet,
\[
\boxed{
P_{e,k}\cos\theta_m+Q_{e,k}\sin\theta_m
\le
S_e^{max}\cos\left(\frac{\pi}{8}\right)+\xi_{e,k}^{S}.}
\]
Here $\xi_{e,k}^{S}\ge0$ is an operational slack.

Define
\[
J_{\rm op}
=
\sum_k\left[
\sum_{i,\phi}
(\xi_{i,\phi,k}^{V,+}+\xi_{i,\phi,k}^{V,-})
+
\sum_e\xi_{e,k}^{S}
\right],
\]
\[
J_{\rm service}
=
\sum_{k=0}^{H-1}\Delta t
\sum_\ell
w_\ell(1-\alpha_{\ell,k})\hat P_{\ell,k}^{L},
\]
and use switching auxiliaries
\[
s_{e,k}\ge z_{e,k}-z_{e,k-1},\qquad
s_{e,k}\ge z_{e,k-1}-z_{e,k}
\]
with
\[
J_{\rm switch}=\sum_{k,e}s_{e,k}.
\]

An optional final voltage-quality tie-break may be retained. The frozen hierarchy is
\[
\boxed{
\lexmin\left(
J_{\rm op},
J_{\rm service},
J_{\rm switch},
J_V^{dev}\ {\rm(optional)}
\right).}
\]

Only operational constraints receive slacks. Fault safety, valid action semantics,
source connectivity, and radiality remain hard constraints.

LinDistFlow is the MPC prediction model; AltDSS is the authoritative nonlinear
unbalanced plant. Therefore
\[
\boxed{
\text{LinDistFlow feasibility}\neq
\text{AltDSS realized feasibility/convergence}.}
\]
MPC solver status and plant convergence are separate diagnostics.

\textbf{REUSE}: dynamic/receding-horizon restoration and updated disturbance/forecast
use from~\cite{Liu2020}.

\textbf{ADAPT}: utility-source restoration; no DER dispatch; explicit priority-weighted
load service; multiphase LinDistFlow MILP; eight-sided thermal polygon; operational
slacks only; full-vector action; explicit plant-model mismatch accounting.

\textbf{NEW software abstraction}: controller-side model compiler, typed MPC
context/output, Gurobi formulation layer, diagnostics, and provenance.

\section{Block 8: masked-DQN RL controller}
Denote the RL controller by $\pi_{P2}^{RL}$.

Controller-side orchestration constructs a circuit-specific finite topology catalogue
\[
\boxed{
\mathcal C=\{\mathbf z^{(1)},\dots,\mathbf z^{(N_C)}\}}
\]
instead of using the monolithic $2^{N_{sw}}$ action space.

At time $t$, a state-dependent structural mask
\[
\boxed{m_t(c)\in\{0,1\}}
\]
excludes catalogue actions that violate hard structural logic such as use of faulted
equipment, radiality, or required source connectivity. The DQN selects
\[
\boxed{
c_t^\star=
\argmax_{c:m_t(c)=1}Q_\theta(s_t,c).}
\]

The RL state contains current common observation and the same load forecast used by MPC:
\[
\boxed{s_t^{RL}=[o_t,\hat{\mathbf D}_{t:t+H-1}].}
\]
Current information includes health/fault status, realized topology, source status,
load $P/Q$, realized voltage/loading, and priorities. Continuous variables are normalized.

The exact origin of the physical action is:
\[
\boxed{
c_t^\star\longrightarrow\mathbf z_t^{req}}
\]
by deterministic catalogue lookup, then
\[
\boxed{
\mathbf z_t^{req}+x_t+\mathcal E_{\rm eval}^{ctrl}
\longrightarrow
\boldsymbol\alpha_{\mathcal L^{curt},t}^{req}}
\]
using the same deterministic priority-aware curtailment logic already frozen for the
graph/reference controller. The baseline RL controller does not solve a separate OPF
merely to compute $\boldsymbol\alpha_t$.

Thus
\[
\boxed{
a_t^{RL}
=
(\mathbf z_t^{req},\boldsymbol\alpha_{\mathcal L^{curt},t}^{req}).}
\]

The authoritative transition is
\[
s_t\xrightarrow{a_t}\text{OpenDSSSim/AltDSS}\rightarrow s_{t+1}.
\]
Replay stores conventional transitions $(s_t,a_t,r_t,s_{t+1},d_t)$. The state may
contain previous realized action/state information needed for Markov sufficiency.

Define normalized priority-weighted unserved service
\[
\boxed{
J_{service,t}^{RL}
=
\frac{\sum_\ell w_\ell(1-\alpha_{\ell,t}^{real})P_{\ell,t}^{L}}
{\sum_\ell w_\ell P_{\ell,t}^{L}}.}
\]

Define normalized switching effort
\[
\boxed{
J_{switch,t}^{RL}
=
\frac{1}{N_{sw}}
\sum_e|z_{e,t}^{real}-z_{e,t-1}^{real}|.}
\]

Realized operational violation quantities may be defined by
\[
\xi_{i,\phi,t}^{V,-}=\max(0,V_i^{min}-|V_{i,\phi,t}|),
\]
\[
\xi_{i,\phi,t}^{V,+}=\max(0,|V_{i,\phi,t}|-V_i^{max}),
\]
\[
\xi_{e,\phi,t}^{I}=\max(0,|I_{e,\phi,t}|-I_{e,\phi}^{max}).
\]

Use documented normalization weights $\omega_V,\omega_I$ to form
\[
J_{op,t}^{RL}=\omega_VJ_{V,t}^{viol}+\omega_IJ_{I,t}^{viol}.
\]

The scalar reward is
\[
\boxed{
r_t=
-\left[
\lambda_{op}J_{op,t}^{RL}
+\lambda_{service}J_{service,t}^{RL}
+\lambda_{sw}J_{switch,t}^{RL}
+\lambda_VJ_{V,t}^{dev}
\right].}
\]
The intended ordering is
\[
\boxed{
J_{op}\succ J_{service}\succ J_{switch}\succ J_V^{dev}\ {\rm(optional)}.}
\]
The $\lambda$ reward coefficients are distinct from service-priority weights $w_\ell$.

The DQN action value is
\[
Q^\pi(s_t,c_t)
=
\mathbb E_\pi\left[
\sum_{j=0}^{\infty}\gamma^jr_{t+j}
\mid s_t,c_t
\right].
\]
With target network $\theta^-$ and replay,
\[
\boxed{
y_t^{DQN}
=
r_t+\gamma(1-d_t)
\max_{c':m_{t+1}(c')=1}Q_{\theta^-}(s_{t+1},c')}
\]
and
\[
\boxed{
L(\theta)
=
\mathbb E[(y_t^{DQN}-Q_\theta(s_t,c_t))^2].}
\]

Training uses $\epsilon$-greedy exploration over admissible actions only. Official TEST
uses $\epsilon=0$.

The backend is frozen as
\[
\boxed{\text{Stable-Baselines3 infrastructure + custom masked-DQN logic}.}
\]
The mask must be respected in both action selection and Bellman target maximization.

\textbf{REUSE}: deep-Q topology selection, online reconfiguration concept, and
switching-frequency concern from~\cite{Oh2020}.

\textbf{ADAPT}: post-fault restoration; priority-weighted service; common MPC/RL
forecast; controller-generated topology catalogue; state-dependent structural mask;
reused rule-based curtailment decoder; realized operational reward; thin plant boundary.

\textbf{NEW software abstraction}: generalized catalogue generator, mask interface,
topology-index mapping, SB3 masked-DQN extension, diagnostics, and provenance.

\section{Block 9: common information and fair-comparison contract}
At a given experiment time, all controllers receive the same realized current physical
information:
\[
o_t^{P2}
=
[
\mathbf h_t,
\mathbf z_{t-1}^{real},
\boldsymbol\sigma_t,
\mathbf P_t^L,
\mathbf Q_t^L,
\mathbf V_t^{real},
\mathbf I_t^{real},
\mathbf w,
\text{time metadata}
].
\]

The exact same feeder, fault realization, initial topology, load realization, sources,
controllable switches, curtailable channels, priorities, bounds, and time axis are used.

MPC and RL receive the identical load forecast:
\[
\boxed{\hat{\mathbf D}_t^{MPC}=\hat{\mathbf D}_t^{RL}.}
\]
The graph/reference controller intentionally uses current information only.

No controller receives realized future loads, realized future voltage/current, or future
fault information not already declared by the scenario.

For official scenario $j$,
\[
\Omega_j
=
\{\text{feeder},\mathcal F_j,\mathbf z_{0,j},
\mathbf P_{0:T,j}^{L},\mathbf Q_{0:T,j}^{L},
\mathbf w_j,\boldsymbol\sigma_{0:T,j},\Delta t,T\},
\]
and
\[
\boxed{\Omega_j^{GR}=\Omega_j^{MPC}=\Omega_j^{RL}.}
\]

Any candidate electrical evaluator used by Reference or the RL curtailment decoder begins
from the current plant state $x_t$. It is a controller-side evaluator, not a future-state
oracle and not an OpenDSSSim service.

Internal MPC objective values and RL reward are not directly compared. Common plant-level
metrics include $\rho^{kW}$, $\rho^{restore}$, $E^U$, $E^{WU}$, $N^{sw}$, voltage
violation magnitude/duration, thermal violation magnitude/duration, realized voltage
extrema, maximum branch loading, convergence diagnostic, controller computation time,
and restoration trajectory/final state.

Intentional algorithmic differences are:
\begin{center}
\begin{tabular}{p{0.23\linewidth}p{0.67\linewidth}}
\toprule
Controller & Internal reasoning\\
\midrule
Graph/reference & Graph search + current-state controller-side electrical evaluation +
deterministic priority-aware curtailment.\\
MPC & MILP + multiphase LinDistFlow + load forecast + joint topology/curtailment decision.\\
RL & SB3 masked DQN + topology catalogue + load forecast + reused deterministic
curtailment decoder.\\
\bottomrule
\end{tabular}
\end{center}

\section{Block 10: controller-side capabilities and orchestration boundary}
OpenDSSSim is a thin authoritative plant interface. At runtime it exposes realized plant
state/health required by the observation/action contract, accepts $a_t^{req}$, performs
the physical transition, and reports realized outcomes/history.

OpenDSSSim does \emph{not} provide controller intelligence, topology search, source
identification, LinDistFlow compilation, candidate evaluation, forecast generation,
curtailment logic, or RL catalogue generation.

Controller-side/orchestration code may read the static OpenDSS circuit folder/files to
derive what the controllers need, including graph topology, fixed/switchable edges,
candidate source buses, phase connectivity, impedances, ratings, load locations,
topology catalogue, and LinDistFlow matrices.

Source identification is controller-side. Curtailability is controller/problem-side.
Operational limits are controller-side problem data derived from justified circuit
metadata or explicit experiment configuration; missing limits are not silently invented
as native feeder facts.

All computational services are controller/orchestration side:
\begin{itemize}
\item topology parser/model compiler;
\item source identification;
\item curtailability declaration;
\item operational-limit construction;
\item current-state LinDistFlow evaluator;
\item forecast provider;
\item MPC LinDistFlow compiler;
\item RL topology catalogue/mask generator.
\end{itemize}

A controller-independent orchestration/problem description may bind
\[
\mathcal P_{P2}
=
\{\text{state schema},\text{action schema},\text{disturbances},
\text{constraints},\text{resources},\text{time semantics}\}.
\]
It contains no graph-algorithm details, Gurobi tuning, or DQN architecture.

Controllers return
\[
\boxed{O_t^{ctrl}=(a_t^{req},d_t,p_t)}
\]
where $d_t$ is diagnostics and $p_t$ is provenance. Only $a_t^{req}$ is required for
physical application.

The orchestration layer lies between controller and plant:
\[
\boxed{
\text{controller}
\rightarrow
\text{orchestration}
\rightarrow
a_t^{req}
\rightarrow
\texttt{env.step}(a_t^{req})
\rightarrow
\text{OpenDSSSim/AltDSS}.}
\]

The orchestration layer routes current state, static models, forecasts, controller
configuration, diagnostics, and action translation. It does not repair the scientific
decision.

\[
\boxed{\texttt{env.step}(a_t^{req})}
\]
is the authoritative physical mutation boundary. Controller code never imports or
manipulates AltDSS to enact actions.

OpenDSSSim records
\[
a_t^{req}\rightarrow a_t^{applied}\rightarrow a_t^{realized}.
\]
The applied action is not a scientifically optimized replacement. The plant may reject
fundamentally inadmissible actuator commands but does not substitute a different
restoration choice.

Missing required static/runtime capabilities produce an explicit machine-readable
\texttt{UNSUPPORTED} result. Support may be controller-specific. Generalized controller
source contains no hard-coded IEEE-123 object names.

\section{Block 11: training, tuning, full-horizon evaluation, and reproducibility}
Every official run has a scenario identity
\[
\Omega_j
=
(\text{feeder version},F_j,L_j,\mathbf z_0,\mathbf w,\boldsymbol\sigma,\Delta t,T).
\]

Only RL has explicit data partitions in this version:
\[
\boxed{
\mathcal D^{RL}
=
\mathcal D^{train}\,\dot\cup\,\mathcal D^{test}.}
\]
There is no VALIDATION partition and no HPO campaign.

Reference and MPC have no train/validation/test fitting process. Their mathematical and
numerical configurations are fixed a priori before official evaluation. For fair
comparison they are run on the same official TEST load trajectories used to evaluate the
frozen RL policy.

Let
\[
\mathcal F=\{F_1,\dots,F_M\}
\]
be a small predetermined set of P2 fault scenarios. For every reported case,
Reference, MPC, and RL use the identical fault condition.

The baseline RL training approach keeps a selected fault scenario fixed and varies TRAIN
load trajectories:
\[
(F_j,L_1^{train}),(F_j,L_2^{train}),\dots.
\]
A scenario-specific RL policy $\pi_{RL}^{(j)}$ is permitted. Training one policy across
multiple faults is a separate experiment.

There is no HPO in this P2 version. RL uses one documented predeclared DQN
configuration. MPC formulation/horizon/tolerances are fixed experiment configuration,
not fitted on TEST results.

Use
\[
\boxed{N_{\rm seed}\ge1.}
\]

Record RL settings including $\gamma$, learning rate, replay size, batch size,
$\epsilon$ schedule, target-network update, architecture, and total training budget.
The frozen checkpoint is chosen by a predeclared TRAIN-only rule, with the final
checkpoint after a fixed training budget as the default. TEST is never used for
checkpoint selection or retuning.

Reward coefficients are fixed before TEST and preserve the intended ordering
\[
J_{op}\succ J_{service}\succ J_{switch}\succ J_V^{dev}\ {\rm(optional)}.
\]

Official evaluation runs all controllers closed loop over the entire common horizon:
\[
\boxed{t=0,\Delta t,\dots,T_{\rm eval}.}
\]
Evaluation does not stop merely because a temporarily good restoration state has been
reached. Loads may continue to change. If the current topology remains best, the
controller may request the same action again, yielding zero switching effort for that
step.

Every P1 and P2 controller instance retains the most recent successfully computed control
command:
\[
\boxed{a_{\rm last}^{success}.}
\]
If the current control computation cannot produce a usable solution under the
controller's declared success rule,
\[
\boxed{a_t^{req}=a_{\rm last}^{success}.}
\]
This is controller memory, not environment correction. Initialize the memory from the
current realized/hold control state so fallback is defined from episode start. Record
\[
f_t^{fallback}\in\{0,1\}.
\]

No silent retry-until-success behavior is permitted. Record controller statuses such as
MPC solver state, reference candidate status, and RL admissible-action status.

Separate controller computation time from plant time:
\[
T_t^{ctrl}\neq T_t^{plant}.
\]

Every official run records feeder/source hash, experiment-config hash, fault ID, load
trajectory ID, controller family/version, controller-config hash, random seed where
applicable, software/environment versions, MPC solver version/settings, and RL
checkpoint hash where applicable.

\section{Frozen cross-controller implementation mapping}
The implementation-facing contract is
\[
\boxed{
\text{current plant state/health}
+
\text{controller-side static model/config}
\rightarrow
\text{controller}
\rightarrow
a_t^{req}
\rightarrow
\text{orchestration}
\rightarrow
\texttt{env.step}.}
\]

All P2 controller families return the same physical action schema:
\[
a_t^{req}
=
(\mathbf z_t^{req},\boldsymbol\alpha_{\mathcal L^{curt},t}^{req}).
\]

The intended controller package organization is:
\begin{verbatim}
src/opendsssim/controllers/
    common/
        base.py
        contracts.py
        capability_requirements.py
        registry.py
        results.py

    restoration/
        reference/
            controller.py
            graph_restoration.py
            config.py

        optimization/
            controller.py
            formulation.py
            config.py

        rl/
            controller.py
            training.py
            policy.py
            config.py
\end{verbatim}

P1 and P2 share these runtime invariants:
\begin{itemize}
\item controllers never mutate AltDSS directly;
\item \texttt{env.step} is the authoritative mutation boundary;
\item requested/applied/realized history is retained;
\item controller intelligence stays out of the environment;
\item every controller retains $a_{\rm last}^{success}$;
\item experiment-specific object IDs live in configuration, not generalized source.
\end{itemize}

\section{Frozen items and experiment-freeze items}
This document freezes the common P2 problem, graph/health/source semantics, priority
objective, full-vector action, controller responsibility for radiality/source connectivity,
operational-slack-only rule, graph/reference formulation, multiphase LinDistFlow MILP
MPC, conservative eight-sided thermal polygon, SB3 masked-DQN RL, fair-information
contract, thin OpenDSSSim boundary, controller-side orchestration services, RL
TRAIN/TEST-only no-HPO protocol, full-horizon evaluation, and last-successful-control memory/fallback.

The following remain deferred to the P2 IEEE-123 experiment freeze:
\begin{itemize}
\item exact P2 working-feeder copy and switch/tie modifications;
\item exact fault scenarios;
\item exact load-priority assignments;
\item exact curtailable-load set;
\item exact timestep from $\{5,15,30,60\}$ min;
\item MPC horizon $H$;
\item exact numerical operating bounds where configuration-supplied;
\item exact operational-slack normalization;
\item exact RL architecture/hyperparameter/reward coefficient values;
\item exact RL seed count, subject to $N_{\rm seed}\ge1$.
\end{itemize}

No deferred choice may silently alter the frozen mathematical problem.

\section{P2 freeze statement}
Blocks 1--11 are frozen for implementation. Controller classes and controller-side
orchestration/model services may now be developed against this contract. Any
implementation change that alters the common physical problem, action meaning,
structural constraints, objective hierarchy, plant/controller boundary, or
fair-information contract requires explicit science-contract change control.

\bibliographystyle{plain}
\bibliography{P2_Reconfiguration_Restoration_References}
\end{document}
